\section{Design methodology and verification setup}
\label{sec:method}

\subsection{Sizing against the real load}
\label{sec:loading}

The decisive step of the design was recognizing what a delay stage actually
drives inside the grid. A $\ttau{2}$ tap feeds the latch inputs of an entire
row, a $\ttau{1}$ tap those of an entire column, and both feed the next delay
stage, whose input capacitance scales with the stage sizing itself. After
fan-out equalization the equivalent load is about twenty unit-inverter inputs
per tap.

Stage sizing derived from a standalone sweep with an idealized load predicted
$\tlsb=\SI{16.6}{\pico\second}$; the same cells assembled into a row with raw
tap loading measured \SI{28.6}{\pico\second} (Section~\ref{sec:1dvs2d}),
nearly a factor of two lost to loading. All final sizing was therefore done
against a replica load of ten latch inputs plus one downstream stage, swept
in the slowest corner (SS) so that the slow line stays reachable everywhere.
The same analysis moved the architecture itself: an earlier $k=4$ operating
point ($\ttau{1}/\ttau{2}=60/\SI{45}{\pico\second}$) required stage delays
below the achievable window and was abandoned in favor of $k=6$.

\subsection{Per-corner trimming}
\label{sec:calmethod}

Process corners shift both line delays together, but not perfectly, so $\tlsb$
breathes with the corner. The MOSCAP banks of Section~\ref{sec:delaycell}
absorb this. All 64 on/off combinations of the six bank gates were simulated
per corner, measuring $\ttau{1}$ and $\ttau{2}$ directly on the tap nets; the
configuration with $\tlsb$ closest to \SI{15}{\pico\second} was selected and
validated end-to-end with a full staircase sweep (Appendix~\ref{app:trim}).
Trimming changes gate voltages only, so the transistor inventory and $\sumw$
are identical in every corner.

\subsection{Testbench and measurement method}
\label{sec:testbench}

All system-level results come from the unmodified course testbench
\texttt{TB\_TDC} \cite{ee4615tb}: stimulus generation in \texttt{TDC\_in},
the DUT wrapper, and \texttt{TDC\_out} loading every thermometer output with
two $2\times$ inverters and \SI{1}{\femto\farad}. The DUT receives its supply
through a dedicated rail with a series ammeter, separate from the stimulus
rail, so the measured energy, $E=\Vdd\int |i|\,dt$ over each
\SI{5}{\nano\second} conversion window, contains only the converter.

OCEAN automates the characterization. Per corner, the input delay is swept in
\SI{1.5}{\pico\second} steps (one tenth of the LSB) across the full range,
and the scripts extract the LSB (mean step width), DNL, INL, offset, energy,
ENOB and the sign-off checks: no missing codes, monotonicity, and an intact
thermometer code at every point. The course script reports INL as the
half-range value $(\mathrm{INL}_{\max}-\mathrm{INL}_{\min})/2$, while the
plots show the endpoint-referenced per-code curve; we quote the course
convention for sign-off and state both where they differ.
